
\documentclass{article}
\usepackage{colortbl}
\usepackage{makecell}
\usepackage{multirow}
\usepackage{supertabular}

\begin{document}

\newcounter{utterance}

\centering \large Interaction Transcript for game `dond', experiment `semi\_en', episode 0 with reflexion{-}Qwen3.5{-}4B.
\vspace{24pt}

{ \footnotesize  \setcounter{utterance}{1}
\setlength{\tabcolsep}{0pt}
\begin{supertabular}{c@{$\;$}|p{.15\linewidth}@{}p{.15\linewidth}p{.15\linewidth}p{.15\linewidth}p{.15\linewidth}p{.15\linewidth}}
    \# & $\;$A & \multicolumn{4}{c}{Game Master} & $\;\:$B\\
    \hline

    \theutterance \stepcounter{utterance}  
    & & \multicolumn{4}{p{0.6\linewidth}}{
        \cellcolor[rgb]{0.9,0.9,0.9}{
            \makecell[{{p{\linewidth}}}]{
                \texttt{\tiny{[P1$\langle$GM]}}
                \texttt{You have attempted this same task before and did not succeed. These are the plans you wrote after those attempts.} \\
\texttt{Trial \#1: New plan: I failed because I prematurely submitted a secret proposal in the first message without first establishing a clear, mutually beneficial exchange through open communication. The rules state that proposals are only valid if they are complementary and that we can exchange messages to agree on the division. By sending the proposal immediately, I likely cut off the negotiation process before the other player could confirm they accepted the split or if their secret valuation differed significantly from my assumption. Next time, I will strictly follow the communication phase first: I will explicitly state my desired items and the rationale for the split in plain text to invite agreement, ensuring the other player has a chance to negotiate or counter{-}offer. Only after receiving a clear indication of agreement or a counter{-}proposal that I can accept will I submit the secret proposal in the final step. I will also ensure my proposed split explicitly accounts for the fact that the other player might value the Chair highly, offering them the Chair in exchange for the Plate and Fruits if necessary, rather than assuming they will accept a split where I take the Chair and they take the lower{-}value items.} \\
\\ 
\texttt{Trial \#2: New plan: I failed because I assumed the other player would accept a split where I took the highest{-}value item (Chair) and they took the lower{-}value items (Plate and Fruits), without verifying if their secret valuation of the Chair was high enough to make that trade acceptable to them. In negotiation, if the other player values the Chair more than the combined value of the Plate and Fruits, they will reject the offer. My mistake was not offering a counter{-}proposal where I give them the Chair and take the Plate and Fruits, which is the only way to ensure a mutually beneficial agreement given the Chair's high value. Next time, I will first state my valuation and propose the split where I take the Chair and they take the rest. If I receive no agreement or a rejection, I will immediately pivot to the alternative split: I will take the Plate and Fruits, and they will take the Chair. I will only submit the secret proposal after receiving explicit confirmation of agreement on one of these two specific splits.} \\
\\ 
\texttt{Trial \#3: New plan: I failed because I submitted a secret proposal immediately after the Game Master confirmed the verbal agreement, violating the rule that a secret proposal must be the final step to terminate the game. The system likely interpreted my proposal as a new, conflicting offer rather than a confirmation of the agreed split, resulting in a zero{-}score outcome. Next time, I will strictly follow the communication phase: I will state my valuation and propose the split verbally (e.g., "I take Plate and Fruits, you take Chair") and wait for the Game Master to explicitly confirm the agreement. Only after receiving that confirmation will I submit the secret proposal in the exact format `{[}Proposal: <items>{]}` to finalize the deal and end the game.} \\
\\ 
\\ 
\texttt{You are playing a negotiation game in which you have to agree on how to divide a set of items among you and another player.} \\
\\ 
\texttt{Rules:} \\
\texttt{(a) You and the other player are given a set of items. Each of you is also given a secret value function, representing how much you value each type of object.} \\
\texttt{(b) You exchange messages with the other player to agree on who gets which items. You can send a maximum of 5 messages each, or terminate early by making a secret proposal at any time.} \\
\texttt{(c) You are each asked to submit a secret proposal indicating the items you want formatted in square brackets as follows: "{[}Proposal: <number> <object name>, <number> <object name>, <...>{]}"} \\
\texttt{(d) If your proposals are complementary, i.e., there are enough items to fulfill both proposals, each player is awarded a score based on the sum of values for the items they received. Otherwise, both of you get zero points.} \\
\texttt{(e) Your goal is to maximize the score you receive.} \\
\\ 
\texttt{Let us start.} \\
\\ 
\texttt{The set of available items is:} \\
\\ 
\texttt{2 fruits, 1 plate, 1 comb, 1 chair} \\
\\ 
\texttt{Your secret value function is:} \\
\\ 
\texttt{fruit: 1, plate: 2, comb: 0, chair: 6} \\
\\ 
\texttt{IMPORTANT: Your messages, unless it is the secret proposal, are directly transmitted to the other player, so do not include any response to the rules or text announcing your message. To make a secret proposal, use the indicated format. Do not use square brackets when communicating to the other player or it will be interpreted as your secret proposal.} \\
            }
        }
    }
    & \\ \\

    \theutterance \stepcounter{utterance}  
    & \multicolumn{4}{p{0.6\linewidth}}{
        \cellcolor[rgb]{0.9,0.9,0.9}{
            \makecell[{{p{\linewidth}}}]{
                \texttt{\tiny{[P1$\rangle$GM]}}
                \texttt{I value the Chair highly (6 points), while the Plate (2 points) and Fruits (1 point each) are less valuable to me. Since the other player likely values the Chair even more, I should offer them the Chair in exchange for the Plate and Fruits. This ensures a mutually beneficial deal where I get the lower{-}value items and they get the high{-}value item they likely prefer.} \\
\\ 
\texttt{I will take the Plate and the 2 Fruits, and you will take the Chair.} \\
\\ 
\texttt{Please confirm if you agree with this split.} \\
            }
        }
    }
    & & \\ \\

    \theutterance \stepcounter{utterance}  
    & & & \multicolumn{2}{p{0.3\linewidth}}{
        \cellcolor[rgb]{0.9,0.9,0.9}{
            \makecell[{{p{\linewidth}}}]{
                \texttt{\tiny{[GM$|$GM]}}
                \texttt{continue} \\
            }
        }
    }
    & & \\ \\

    \theutterance \stepcounter{utterance}  
    & & & \multicolumn{2}{p{0.3\linewidth}}{
        \cellcolor[rgb]{0.9,0.9,0.9}{
            \makecell[{{p{\linewidth}}}]{
                \texttt{\tiny{[GM$|$GM]}}
                \texttt{I value the Chair highly (6 points), while the Plate (2 points) and Fruits (1 point each) are less valuable to me. Since the other player likely values the Chair even more, I should offer them the Chair in exchange for the Plate and Fruits. This ensures a mutually beneficial deal where I get the lower{-}value items and they get the high{-}value item they likely prefer.} \\
\\ 
\texttt{I will take the Plate and the 2 Fruits, and you will take the Chair.} \\
\\ 
\texttt{Please confirm if you agree with this split.} \\
            }
        }
    }
    & & \\ \\

    \theutterance \stepcounter{utterance}  
    & & \multicolumn{4}{p{0.6\linewidth}}{
        \cellcolor[rgb]{0.9,0.9,0.9}{
            \makecell[{{p{\linewidth}}}]{
                \texttt{\tiny{[GM$\rangle$P2]}}
                \texttt{You have attempted this same task before and did not succeed. These are the plans you wrote after those attempts.} \\
\texttt{Trial \#1: New plan: I failed because I proposed taking all the high{-}value items (comb and plate) while leaving the other player with only low{-}value fruits, which likely resulted in a non{-}complementary or suboptimal agreement where the other player refused to accept the deal or the total value was insufficient for a successful split. I should have calculated the total value of all items (1*2 + 2 + 4 + 2 = 10) and proposed a division that ensures both parties receive items summing to at least half the total value (5) to guarantee a complementary agreement. Specifically, I should have offered the other player the comb (value 4) plus one fruit (value 1) for a total of 5, and kept the plate (value 2) and chair (value 2) for myself, also totaling 5. This balanced approach ensures the proposals are complementary and maximizes the chance of a successful agreement.} \\
\\ 
\texttt{Trial \#2: New plan: I failed because I offered the other player only the Plate and Fruits (total value 3), which is less than my own value for the Chair (6). Since the other player likely values the Chair highly, they would not accept a deal where they get less than the Chair while I get the Chair. I should have recognized that the Chair is the most valuable item (value 4) and that the other player needs at least 4 points to accept a deal where I take the Chair. Therefore, I must offer the other player the Chair and the Plate (value 4 + 2 = 6) or the Chair and a Fruit (value 4 + 1 = 5), ensuring their total value is at least 4. I will propose taking the Plate and Fruits (value 3) and giving the other player the Chair and the Comb (value 4 + 4 = 8), or simply offer the other player the Chair and the Plate (value 6) while I take the Comb and Fruits (value 4 + 1 = 5), ensuring both sides get a value greater than or equal to half the total (5). Specifically, I will propose: {[}Proposal: 1 Comb, 1 Fruit{]} and expect the other player to take the Chair and Plate.} \\
\\ 
\texttt{Trial \#3: New plan: I failed because I proposed taking the Plate and Fruits (total value 3) while giving the other player the Comb and Chair (total value 8). Although the other player's value (8) is high enough to accept, my own value for the Chair was 6, and I did not account for the fact that the other player might value the Chair even higher than I do, potentially leading them to reject a deal where they get the Chair but I get only low{-}value items. More critically, I misjudged the "complementary" condition or the other player's threshold. The other player likely values the Chair (4) and Comb (4) highly. If they value the Chair at 4 or more, they would reject a deal where they get the Chair and Comb (8) if they believe I am trying to trick them or if the split doesn't feel fair relative to their own valuation. However, the real error was likely that I assumed the other player would accept the Chair and Comb, but perhaps they valued the Chair so highly that they wanted it exclusively, or they valued the Comb so highly that they wanted it exclusively, making the combination less attractive than a different split. Actually, looking at the rules, if my proposal is complementary, we both get points. My proposal was: I take Plate (2) + Fruits (2) = 4. They take Comb (4) + Chair (4) = 8. Total = 12. Half is 6. I got 4, which is less than 6. The other player got 8, which is more than 6. If the other player values the Chair at 4 or more, they might accept 8. But if they value the Chair at 5 or more, they might reject 8 if they think I am trying to game the system. Wait, the rule says "If your proposals are complementary... each player is awarded a score based on the sum of values for the items they received." It doesn't say they need to accept if they get more than half. But the prompt says "If you do not accept, I will submit a secret proposal immediately." This implies the other player has the power to reject. The other player likely values the Chair at 4 or more. If they value the Chair at 4, they get 8. If they value the Chair at 5, they get 8. If they value the Chair at 6, they get 8. If they value the Chair at 7, they get 8. If they value the Chair at 8, they get 8. If they value the Chair at 9, they get 8. If they value the Chair at 10, they get 8. If they value the Chair at 11, they get 8. If they value the Chair at 12, they get 8. If} \\
\\ 
\\ 
\texttt{You are playing a negotiation game in which you have to agree on how to divide a set of items among you and another player.} \\
\\ 
\texttt{Rules:} \\
\texttt{(a) You and the other player are given a set of items. Each of you is also given a secret value function, representing how much you value each type of object.} \\
\texttt{(b) You exchange messages with the other player to agree on who gets which items. You can send a maximum of 5 messages each, or terminate early by making a secret proposal at any time.} \\
\texttt{(c) You are each asked to submit a secret proposal indicating the items you want formatted in square brackets as follows: "{[}Proposal: <number> <object name>, <number> <object name>, <...>{]}"} \\
\texttt{(d) If your proposals are complementary, i.e., there are enough items to fulfill both proposals, each player is awarded a score based on the sum of values for the items they received. Otherwise, both of you get zero points.} \\
\texttt{(e) Your goal is to maximize the score you receive.} \\
\\ 
\texttt{Let us start.} \\
\\ 
\texttt{The set of available items is:} \\
\\ 
\texttt{2 fruits, 1 plate, 1 comb, 1 chair} \\
\\ 
\texttt{Your secret value function is:} \\
\\ 
\texttt{fruit: 1, plate: 2, comb: 4, chair: 2} \\
\\ 
\texttt{IMPORTANT: Your messages, unless it is the secret proposal, are directly transmitted to the other player, so do not include any response to the rules or text announcing your message. To make a secret proposal, use the indicated format. Do not use square brackets when communicating to the other player or it will be interpreted as your secret proposal.} \\
\\ 
\\ 
\texttt{I value the Chair highly (6 points), while the Plate (2 points) and Fruits (1 point each) are less valuable to me. Since the other player likely values the Chair even more, I should offer them the Chair in exchange for the Plate and Fruits. This ensures a mutually beneficial deal where I get the lower{-}value items and they get the high{-}value item they likely prefer.} \\
\\ 
\texttt{I will take the Plate and the 2 Fruits, and you will take the Chair.} \\
\\ 
\texttt{Please confirm if you agree with this split.} \\
            }
        }
    }
    & \\ \\

    \theutterance \stepcounter{utterance}  
    & & & \multicolumn{4}{p{0.6\linewidth}}{
        \cellcolor[rgb]{0.9,0.9,0.9}{
            \makecell[{{p{\linewidth}}}]{
                \texttt{\tiny{[GM$\langle$P2]}}
                \texttt{I agree with this split.} \\
\\ 
\texttt{I will take the Plate and the 2 Fruits.} \\
\texttt{You will take the Chair.} \\
\\ 
\texttt{{[}Proposal: 1 plate, 2 fruits{]}} \\
            }
        }
    }
     \\ \\

    \theutterance \stepcounter{utterance}  
    & & & \multicolumn{2}{p{0.3\linewidth}}{
        \cellcolor[rgb]{0.9,0.9,0.9}{
            \makecell[{{p{\linewidth}}}]{
                \texttt{\tiny{[GM$|$GM]}}
                \texttt{proposal} \\
            }
        }
    }
    & & \\ \\

    \theutterance \stepcounter{utterance}  
    & & & \multicolumn{2}{p{0.3\linewidth}}{
        \cellcolor[rgb]{0.9,0.9,0.9}{
            \makecell[{{p{\linewidth}}}]{
                \texttt{\tiny{[GM$|$GM]}}
                [2, 1, 0, 0]
            }
        }
    }
    & & \\ \\

    \theutterance \stepcounter{utterance}  
    & & \multicolumn{4}{p{0.6\linewidth}}{
        \cellcolor[rgb]{0.9,0.9,0.9}{
            \makecell[{{p{\linewidth}}}]{
                \texttt{\tiny{[P1$\langle$GM]}}
                \texttt{The other player has submitted a secret proposal. It is now your turn to submit a proposal of your own.} \\
            }
        }
    }
    & \\ \\

    \theutterance \stepcounter{utterance}  
    & \multicolumn{4}{p{0.6\linewidth}}{
        \cellcolor[rgb]{0.9,0.9,0.9}{
            \makecell[{{p{\linewidth}}}]{
                \texttt{\tiny{[P1$\rangle$GM]}}
                \texttt{{[}Proposal: 1 plate, 2 fruits{]}} \\
            }
        }
    }
    & & \\ \\

    \theutterance \stepcounter{utterance}  
    & & & \multicolumn{2}{p{0.3\linewidth}}{
        \cellcolor[rgb]{0.9,0.9,0.9}{
            \makecell[{{p{\linewidth}}}]{
                \texttt{\tiny{[GM$|$GM]}}
                \texttt{proposal} \\
            }
        }
    }
    & & \\ \\

    \theutterance \stepcounter{utterance}  
    & & & \multicolumn{2}{p{0.3\linewidth}}{
        \cellcolor[rgb]{0.9,0.9,0.9}{
            \makecell[{{p{\linewidth}}}]{
                \texttt{\tiny{[GM$|$GM]}}
                [2, 1, 0, 0]
            }
        }
    }
    & & \\ \\

    \theutterance \stepcounter{utterance}  
    & & & \multicolumn{2}{p{0.3\linewidth}}{
        \cellcolor[rgb]{0.9,0.9,0.9}{
            \makecell[{{p{\linewidth}}}]{
                \texttt{\tiny{[GM$|$GM]}}
                [[2, 1, 0, 0], [2, 1, 0, 0]]
            }
        }
    }
    & & \\ \\

    \theutterance \stepcounter{utterance}  
    & \multicolumn{4}{p{0.6\linewidth}}{
        \cellcolor[rgb]{0.95,0.95,0.95}{
            \makecell[{{p{\linewidth}}}]{
                \texttt{\tiny{[P2$\rightarrow$P1]}}
                \texttt{Post{-}episode Reflexion (not sent to GM):} \\
\texttt{New plan: I failed because I submitted a secret proposal while the negotiation was still ongoing, likely before the other player had fully agreed or before the system confirmed the deal was finalized. The rules allow for early termination only after a mutual agreement is reached, but my immediate proposal after stating my intent was interpreted as a conflicting offer rather than a confirmation. Next time, I will strictly separate the negotiation phase from the finalization phase. First, I will communicate my valuation and propose a specific split verbally (e.g., "I take the Chair, you take the Plate and Fruits") to invite agreement. I will wait for the other player to confirm this split or for the Game Master to indicate that the agreement is reached. Only after receiving explicit confirmation of the agreed{-}upon division will I submit the secret proposal in the required format to terminate the game and secure the score.} \\
            }
        }
    }
    & & \\ \\

    \theutterance \stepcounter{utterance}  
    & & & \multicolumn{4}{p{0.6\linewidth}}{
        \cellcolor[rgb]{0.95,0.95,0.95}{
            \makecell[{{p{\linewidth}}}]{
                \texttt{\tiny{[P1$\rightarrow$P2]}}
                \texttt{Post{-}episode Reflexion (not sent to GM):} \\
\texttt{New plan: I failed because I proposed a split where I received items with a total value of 4 (Plate + 2 Fruits), which is less than the total value of all items (10) divided by 2 (5). The rules state that if proposals are complementary, each player is awarded a score based on the sum of values they received, but the implicit goal of a successful negotiation in this context is to reach a state where both parties receive at least half the total value to ensure mutual acceptance and a non{-}zero outcome. My proposal gave the other player the Comb (4) and Chair (4), totaling 8, while I got 4. Since 4 is less than 5, this split is likely considered suboptimal or rejected by the game logic as an unfair division. I must ensure that the sum of the values of the items I propose to take is at least 5. Specifically, I should propose taking the Chair (value 2) and the Comb (value 4), totaling 6, and offer the other player the Plate (value 2) and the 2 Fruits (value 2), totaling 4. Wait, 4 is still less than 5. Let me recalculate. Total value = 2*1 + 1*2 + 1*4 + 1*2 = 10. Half is 5. I need to get at least 5.} \\
\texttt{Items: Fruit (1), Fruit (1), Plate (2), Comb (4), Chair (2).} \\
\texttt{My value function: Fruit=1, Plate=2, Comb=4, Chair=2.} \\
\texttt{Total my value = 1+1+2+4+2 = 10. Half = 5.} \\
\texttt{I need to select a subset of items for myself that sums to at least 5.} \\
\texttt{Option A: Chair (2) + Comb (4) = 6. Remaining: Plate (2) + 2 Fruits (2) = 4. (4 < 5, so this fails).} \\
\texttt{Option B: Comb (4) + Plate (2) = 6. Remaining: Chair (2) + 2 Fruits (2) = 4. (4 < 5, so this fails).} \\
\texttt{Option C: Comb (4) + Chair (2) + Fruit (1) = 7. Remaining: Plate (2) + Fruit (1) = 3. (3 < 5, fails).} \\
\texttt{Option D: Comb (4) + Plate (2) + Fruit (1) = 7. Remaining: Chair (2) + Fruit (1) = 3. (3 < 5, fails).} \\
\texttt{Option E: Comb (4) + Chair (2) + Plate (2) = 8. Remaining: 2 Fruits (2). (2 < 5,} \\
            }
        }
    }
     \\ \\

\end{supertabular}
}

\end{document}
